% design_sistema.tex — Documento di Design del Sistema

\section{Design del Sistema}

\subsection{Descrizione dell'architettura}

BugBoard26 adotta un'architettura \textbf{client-server a tre livelli}
(Three-Tier Architecture), separando nettamente le responsabilità tra
\emph{presentazione}, \emph{logica di business} e \emph{persistenza dei dati}.

\begin{center}
\begin{tabular}{|l|p{10cm}|}
\hline
\textbf{Livello} & \textbf{Descrizione} \\
\hline
Presentazione (Frontend) & Single Page Application React che comunica con il backend tramite API REST. \\
\hline
Logica di Business (Backend) & Applicazione Spring Boot che espone endpoint RESTful, gestisce autenticazione JWT e implementa le regole di business. \\
\hline
Persistenza (Database) & Database relazionale (H2 in sviluppo, PostgreSQL in produzione) gestito tramite JPA/Hibernate. \\
\hline
\end{tabular}
\end{center}

\subsubsection{Diagramma architetturale}

\begin{verbatim}
  +-----------------+       HTTPS/REST       +-------------------+
  |                 |  <------------------->  |                   |
  |    Frontend     |       /api/*           |     Backend       |
  |    (React SPA)  |                        |   (Spring Boot)   |
  |    Port 8080    |                        |    Port 8081      |
  +-----------------+                        +--------+----------+
        |                                             |
        | (Vite dev proxy)                            | JPA/Hibernate
        | /api -> localhost:8081                       |
        |                                    +--------v----------+
        |                                    |                   |
        |                                    |     Database      |
        |                                    |  H2 / PostgreSQL  |
        |                                    +-------------------+
        |
        v
  +-----------------+
  | Docker Compose  |
  |  (Produzione)   |
  +-----------------+
\end{verbatim}

\subsection{Criteri di design adottati}

\begin{description}
    \item[Separazione delle responsabilità (\emph{Separation of Concerns})]
    Ogni livello ha una responsabilità ben definita: il frontend si occupa
    esclusivamente della presentazione e dell'interazione utente, il backend
    della logica di business e dell'autorizzazione, il database della
    persistenza. Questa separazione favorisce la manutenibilità e la
    testabilità.

    \item[Stateless Authentication]
    L'autenticazione è gestita tramite token JWT, eliminando la necessità
    di sessioni server-side. Ogni richiesta contiene il token necessario
    per l'identificazione, facilitando la scalabilità orizzontale.

    \item[API-First Design]
    Il backend espone un'API RESTful documentata tramite OpenAPI/Swagger
    (SpringDoc). Questo approccio consente lo sviluppo indipendente di
    frontend e backend, facilita l'integrazione con altri client e permette
    il testing automatico delle API.

    \item[Convention over Configuration]
    L'utilizzo di Spring Boot riduce drasticamente la configurazione
    boilerplate grazie alle \emph{auto-configurations} e alle convenzioni
    del framework.

    \item[Containerizzazione]
    L'intero stack è containerizzabile tramite Docker e orchestrabile con
    Docker Compose, garantendo riproducibilità dell'ambiente e facilità di
    deployment.
\end{description}

\subsection{Scelte tecnologiche}

\subsubsection{Frontend}

\begin{center}
\small
\begin{tabular}{|l|l|p{6cm}|}
\hline
\textbf{Tecnologia} & \textbf{Versione} & \textbf{Motivazione} \\
\hline
React & 18.x & Libreria UI component-based matura e con vasto ecosistema; facilita la creazione di SPA reattive. \\
\hline
TypeScript & 5.x & Tipizzazione statica che riduce i bug a runtime e migliora la manutenibilità del codice. \\
\hline
Vite & 6.x & Build tool moderno, significativamente più veloce di Webpack grazie a ESBuild e HMR nativo. \\
\hline
Tailwind CSS & 3.x & Framework CSS utility-first che accelera lo sviluppo della UI e garantisce consistenza visiva. \\
\hline
Radix UI & -- & Componenti accessibili e non stilizzati (\emph{headless}), conformi a WAI-ARIA, personalizzabili con Tailwind. \\
\hline
React Query & 5.x & Gestione dello stato server-side (caching, refetching, invalidation) che semplifica le chiamate API. \\
\hline
React Router & 7.x & Routing dichiarativo per SPA con supporto a lazy loading e route protette. \\
\hline
\end{tabular}
\end{center}

\subsubsection{Backend}

\begin{center}
\small
\begin{tabular}{|l|l|p{6cm}|}
\hline
\textbf{Tecnologia} & \textbf{Versione} & \textbf{Motivazione} \\
\hline
Java & 23 & Linguaggio robusto e type-safe, con record e pattern matching per DTO concisi. \\
\hline
Spring Boot & 3.3.4 & Framework enterprise standard per applicazioni Java; auto-configuration, starter dependencies, embedded server. \\
\hline
Spring Security & -- & Gestione completa di autenticazione e autorizzazione; integrazione nativa con JWT. \\
\hline
Spring Data JPA & -- & Astrazione ORM che elimina il codice boilerplate per l'accesso ai dati; supporto a Specification per query dinamiche. \\
\hline
H2 Database & -- & Database embedded per lo sviluppo locale; nessuna installazione richiesta, avvio immediato. \\
\hline
PostgreSQL & 16 & Database relazionale per produzione; robusto, conforme ACID, supporto nativo a UUID. \\
\hline
JJWT & 0.12.5 & Libreria per la creazione e validazione di token JWT; API fluente e sicura. \\
\hline
SpringDoc & 2.6.0 & Generazione automatica della documentazione OpenAPI 3.0 dalle annotazioni Spring. \\
\hline
Apache POI & 5.3.0 & Generazione di file Excel (.xlsx) per l'export dei dati bug. \\
\hline
\end{tabular}
\end{center}

\subsubsection{Testing}

\begin{center}
\small
\begin{tabular}{|l|p{10cm}|}
\hline
\textbf{Tecnologia} & \textbf{Motivazione} \\
\hline
JUnit 5 & Framework di testing standard per Java; supporto a test parametrizzati, display names, extensions. \\
\hline
Mockito 5 & Libreria di mocking che consente l'isolamento delle dipendenze nei test di unità. \\
\hline
AssertJ & Asserzioni fluenti che migliorano la leggibilità dei test. \\
\hline
\end{tabular}
\end{center}

\subsubsection{DevOps e Deployment}

\begin{center}
\small
\begin{tabular}{|l|p{10cm}|}
\hline
\textbf{Tecnologia} & \textbf{Motivazione} \\
\hline
Docker & Containerizzazione per ambienti riproducibili; Dockerfile multi-stage per build ottimizzati. \\
\hline
Docker Compose & Orchestrazione dei tre servizi (database, backend, frontend) con un solo comando. \\
\hline
Git / GitHub & Version control distribuito; collaborazione, code review, CI/CD integration. \\
\hline
Maven & Build automation per il backend; gestione dipendenze, lifecycle standardizzato. \\
\hline
npm & Package manager per il frontend; gestione dipendenze e script di build. \\
\hline
\end{tabular}
\end{center}

\subsection{Comunicazione tra i livelli}

La comunicazione tra frontend e backend avviene tramite API REST su HTTP/HTTPS:

\begin{itemize}
    \item \textbf{Formato dati}: JSON (Content-Type: \texttt{application/json})
    \item \textbf{Autenticazione}: Token JWT trasportato tramite cookie HttpOnly
          (\texttt{token}) oppure header \texttt{Authorization: Bearer <token>}
    \item \textbf{Proxy in sviluppo}: Vite redirige \texttt{/api/*} a
          \texttt{localhost:8081}, evitando problemi CORS durante lo sviluppo
    \item \textbf{Endpoint principali}:
    \begin{itemize}
        \item \texttt{POST /api/auth/login} — Autenticazione
        \item \texttt{GET/POST /api/bugs} — Listing e creazione bug
        \item \texttt{PATCH /api/bugs/\{id\}} — Modifica bug
        \item \texttt{POST /api/bugs/\{id\}/assign} — Assegnazione (admin)
        \item \texttt{GET /api/admin/metrics} — Metriche (admin)
        \item \texttt{GET /api/export/bugs.csv} — Export CSV
    \end{itemize}
\end{itemize}

\subsection{Sicurezza}

\begin{itemize}
    \item \textbf{Password}: hash BCrypt con salt automatico
    \item \textbf{CSRF}: disabilitato (architettura stateless con JWT)
    \item \textbf{CORS}: configurazione restrittiva con whitelist degli origin consentiti
    \item \textbf{Sessioni}: nessuna sessione server-side (\texttt{SessionCreationPolicy.STATELESS})
    \item \textbf{Autorizzazione}: controllo a livello di servizio basato sul
          ruolo dell'utente (ADMIN, USER, READONLY)
    \item \textbf{Cookie JWT}: attributi \texttt{HttpOnly} e \texttt{SameSite=Lax}
          per mitigare XSS e CSRF
\end{itemize}
